\section{实验总结}
\indent 这次实验围绕 Enigma 密码机的加密实现、密钥恢复和密钥空间分析展开。通过模拟接线板、扰频器与反射器及其反向过程，
我们进一步理解了 Enigma 的加密结构以及转子状态变化对加密结果的影响。在给定参数下，程序将明文 WETTER 加密为 TWNVCX，
与预期结果一致，验证了该配置下加密流程的正确性。

在密钥恢复部分，我们在中速和慢速转子不转动的假设下，利用已知明密文中的特殊的Crib环路关系建立约束，
将转子设置的搜索与接线板映射的恢复分开处理，对密钥空间进行分割。通过多个环路筛选候选，再利用完整明密文关系扩展映射、排除矛盾，最终
得到转子初始位置 \((3,14,7)\) 及部分接线板映射。这一过程使我们认识到，密码分析可以利用密码结构与已知信息对搜索范围进行分割，
而较大的理论密钥空间并不意味着能够抵抗所有分析方法。

在密钥空间与性能分析部分，当从四个转子中选择三个并允许接线板最多使用六条接线时，
总密钥空间约为 \(2^{55.31}\)。根据实验测得的连续加密吞吐量进行粗略估算，遍历完整密钥空间需要约 \(1.75\times10^4\) 年，说明在本实验的实现与现有的硬件条件下，
在个人计算机上直接穷举任务三的完整密钥空间完全不可行。


通过本次实验，我们加深了对于 Enigma 密码机各个部件及其作用的理解，也体会到了程序实现、结果验证与复杂度分析需要相互配合。
后续可通过增加已知明密文补全尚未确定的接线板映射，并在更完整的转子转动条件下检验攻击方法的适用性。